\BourbakiExercisesSection

\begin{enumerate}
\def\labelenumi{\arabic{enumi}.}
\tightlist
\item
  在有限群 G 中，证明一个元素 \(a \in G\) 的共轭元素（第 4 号）的个数等于 a 的正规化子的指数，从而是 G 的阶的一个因子。
\end{enumerate}

¶ 2. *若 G 是一个阶为 n 的有限群，则 G 的自同构个数为 \(\leqslant n^{\log n/\log 2}\) （证明存在 G 的一个生成集 \(\{a_{1}, a_{2}, \ldots, a_{m}\}\) ，使得 \(a_{i}\) 不属于由 \(a_{1}, a_{2}, \ldots, a_{i-1}\) 生成的子群（对 \(2 \leqslant i \leqslant m\) ）；由此推出 \(2^{m} \leqslant n\) ，且 G 的自同构个数为 \(\leqslant n^{m})\cdot*\)

\begin{enumerate}
\def\labelenumi{\arabic{enumi}.}
\setcounter{enumi}{2}
\tightlist
\item
  设 \(\Gamma\) 是一个群 \(G\) 的自同构群， \(\Delta\) 是 \(G\) 的内自同构群；证明 \(\Delta\) 是 \(\Gamma\) 的一个正规子群。为使一个自同构 \(\sigma\) （ \(G\) 的）与 \(G\) 的所有内自同构可交换，必要且充分的是：对所有 \(x \in G\) ， \(x^{-1}\sigma(x)\) 属于 \(G\) 的中心；由此推出，若 \(G\) 的中心约化为单位元，则 \(\Delta\) 在 \(\Gamma\) 中的中心化子也约化为单位元。
\end{enumerate}

¶ 4. (a) 设 G 是一个非交换单群，Γ 是 G 的自同构群，Δ 是 G 的内自同构群（同构于 G）。若 s 是群 Γ 的一个自同构，证明 \(s(\Delta) = \Delta\) （利用上面的习题 3 以及 § 4，第 9 号，命题 15，注意交 \(\Delta \cap s(\Delta)\) 不可能仅由 Γ 的单位元组成）。

\begin{enumerate}
\def\labelenumi{(\alph{enumi})}
\setcounter{enumi}{1}
\item
  证明 \(\Gamma\) 的使 \(\Delta\) 的每个元素都不变的唯一自同构是恒等映射（利用该自同构使 \(\alpha\) 与 \(\sigma \alpha \sigma^{-1}\) （对所有 \(\alpha \in \Delta\) 与 \(\sigma \in \Gamma\) ）都不变这一事实，并使用习题 3）。
\item
  设 \(s\) 是 \(\Gamma\) 的一个自同构， \(\phi\) 是同构 \(x \mapsto \operatorname{Int}(x)\) （G 到 \(\Delta\) 上的）， \(\psi\) 是其逆同构， \(\sigma\) 是 G 的自同构 \(\psi \circ s \circ \phi\) ；证明自同构 \(\xi \mapsto \sigma^{-1}s(\xi)\sigma\) （ \(\Gamma\) 上的）是恒等自同构（利用 (b)，并注意到，对所有 \(x \in G\) ， \(s(\operatorname{Int}(x)) = \operatorname{Int}(\sigma(x)))\) 。由此推出， \(\Gamma\) 的每个自同构都是内自同构。
\end{enumerate}

5. 设 G 是一个群。

\begin{enumerate}
\def\labelenumi{(\alph{enumi})}
\item
  若 H 是 G 的一个指数为有限数 n 的子群，证明 H 的诸共轭子群的交 N 的指数是 \(n!\) 的一个约数（注意 G/N 同构于 \(S_{n}\) 的一个子群）。
\item
  G 称为剩余有限的，如果它的有限指数子群的交是 \(\{e\}\) 。证明这等价于说 G 同构于有限群的一个乘积的一个子群。
\end{enumerate}
% semantic-fragments: legacy-input-seam
\relax{}
\begin{enumerate}
\def\labelenumi{(\alph{enumi})}
\setcounter{enumi}{2}
\item
  假设 G 可以由一个有限集合生成。证明：对每一个整数 n，由 G 的指数为 n 的子群组成的集合 \(P_{n}\) 是有限的。
\item
  在(c)的假设下，设 \(f\) 是 \(G\) 的一个满射自同态。证明：对每个整数 \(n\)，映射 \(H \mapsto f^{-1}(H)\) 是 \(P_n\) 到自身的一个双射（注意这是一个单射）。由此推出，\(f\) 的核包含于 \(G\) 的每一个有限指数子群中。特别地，如果 \(G\) 是剩余有限的，那么 \(f\) 是双射。
\end{enumerate}

\begin{enumerate}
\def\labelenumi{\arabic{enumi}。}
\setcounter{enumi}{5}
\tightlist
\item
  设 H 是群 G 中指数有限的一个子群。假设 G 是 H 的诸共轭子群之并。证明 H = G。（利用习题 5 把它归结为 G 为有限群的情形。在后一情形，注意：
\end{enumerate}

\[
\operatorname{Card} \left(\bigcup_ {x \in G} x H x ^ {- 1}\right) \leqslant 1 + (\operatorname{Card} (H) - 1) \operatorname{Card} (G / H).)
\]

\begin{enumerate}
\def\labelenumi{\arabic{enumi}。}
\setcounter{enumi}{6}
\tightlist
\item
  设 \(p\) 是一个素数。
\end{enumerate}

\begin{enumerate}
\def\labelenumi{(\alph{enumi})}
\item
  证明 \(n^p \equiv n \pmod{p}\) 对所有 \(n \in \mathbf{Z}\) 成立。（先归结为 \(n\) 为正的情形，再对 \(n\) 利用二项式公式进行归纳论证。）
\item
  设 \(x, y\) 是群 \(G\) 的两个元素。假设
\end{enumerate}

\[
y x y ^ {- 1} = x ^ {n}, \quad \text { with } n \in \mathbf {Z}, \quad \text { and   that } \quad x ^ {p} = 1.
\]

证明 \(y^{p}xy^{-p} = x^{n}\)；由此推出 \(y^{p - 1}\) 与 \(x\) 交换。

\begin{enumerate}
\def\labelenumi{(\alph{enumi})}
\setcounter{enumi}{2}
\tightlist
\item
  设 G 是一个群，其所有 \(\neq1\) 的元素的阶均为 p 且彼此共轭。证明 \(\operatorname{Card}(G)\) 等于 1 或 2。（利用 (b) 证明 G 是交换的。）
\end{enumerate}

\begin{enumerate}
\def\labelenumi{\arabic{enumi}。}
\setcounter{enumi}{7}
\tightlist
\item
  设 A 和 B 是有限群 G 的子群， \(N_{A}\) 和 \(N_{B}\) 是 A 和 B 的正规化子， \(v_{A}\) 和 \(v_{B}\) 是 \(N_{A}\) 和 \(N_{B}\) 在 G 中的指数， \(r_{A}\) 是包含 B 的 A 的共轭子群个数， \(r_{B}\) 是包含 A 的 B 的共轭子群个数；证明 \(v_{A}r_{B} = v_{B}r_{A}\)。
\end{enumerate}

设 G 是有限集合 E 的一个置换群；对所有 \(s \in G\)，设 \(\chi(s)\) 表示 \(s\) 的不动点集合的基数。

\begin{enumerate}
\def\labelenumi{(\alph{enumi})}
\item
  证明 \(\sum_{s\in \mathbf{G}}\chi (s) = \mathrm{N}t\)，其中 \(\mathbf{N} = \operatorname {Card}(\mathbf{G})\)，而 \(t\) 是 \(\mathbf{G}\) 在 \(\mathbf{E}\) 中轨道的个数（以两种方式计算有序对 \((s,x)\in \mathbf{G}\times \mathbf{E}\) 中满足 \(s(x) = x\) 者的个数）。
\item
  假设对所有 \(a \in \mathbf{E}\)，记 \(\mathrm{H}_a\) 为 \(a\) 的稳定化子，它不约化为 \(e\)。证明：如果 \(\chi(s)\) 与 \(s\) 无关（对 \(s \neq e\) 而言）且等于某个整数 \(k\)，则
\end{enumerate}

\[
k \leqslant t \leqslant 2 k.
\]

（注意 \(\operatorname{Card}(\mathbf{E}) \leqslant \mathbf{N}k\)。）

在 k = 2 的特殊情形下，求这些群 \(H_{a}\) 的可能的阶；证明，若 t = 3，则对三个轨道中两个轨道的元素而言， \(H_{a}\) 的阶不可能 \(\geqslant 3\)，除非 N 取 12、24 或 60 诸值之一。

\begin{enumerate}
\def\labelenumi{(\alph{enumi})}
\setcounter{enumi}{2}
\tightlist
\item
  假设 G 传递地作用于 E 上，并设 H 表示 E 中某个元素的稳定化子；证明 \(\sum_{s\in G}\chi(s)^{2}\) 等于 N. \(t_{H}\)，其中 \(t_{H}\) 是 H 在 E 中的轨道个数。（对每个 \((s,u)\)，记 \(\chi(s,u)=\chi(s)\)。以两种方式求和 \(\sum_{(s,u)\in R}\chi(s,u)\)，其中 R 是由有序对 \((s,u)\) 组成的集合，这些有序对满足 \(usu^{-1}\in H.\)）
\end{enumerate}

\begin{enumerate}
\def\labelenumi{\arabic{enumi}。}
\setcounter{enumi}{9}
\tightlist
\item
  \begin{enumerate}
  \def\labelenumii{(\alph{enumii})}
  \tightlist
  \item
    证明 \(\tau_{12}\tau_{34},\tau_{13}\tau_{24},\tau_{14}\tau_{23}\) 这些交错群 \(\mathfrak{A}_4\) 的元素与单位元一起构成 \(\mathfrak{A}_4\) 的一个交换子群 H。证明 H 在 \(\mathfrak{A}_4\) 中和 \(\mathfrak{S}_4\) 中都是正规的，并且 \(\mathfrak{A}_4 / \mathrm{H}\) 是 3 阶循环群。
  \end{enumerate}
\end{enumerate}

\begin{enumerate}
\def\labelenumi{(\alph{enumi})}
\setcounter{enumi}{1}
\item
  证明 H 在 \(S_{4}\) 中的中心化子等于 H。由此推出，映射 \(s \mapsto (h \mapsto shs^{-1})\) 把 \(S_{4}\) 映入 \(\operatorname{Aut}(H)\)，并在过渡到商群时定义了 \(S_{4}/H\) 到 \(\operatorname{Aut}(H)\) 上的一个同构，且后一群同构于 \(S_{3}\)。
\item
  设 \(\mathbf{K}\) 是 \(\mathbf{H}\) 的一个 2 阶子群。证明 \(\mathbf{K}\) 在 \(\mathfrak{A}_4\) 中不正规。
\end{enumerate}

\begin{enumerate}
\def\labelenumi{\arabic{enumi}。}
\setcounter{enumi}{10}
\tightlist
\item
  设 E 为有限集， \(\zeta \in \mathfrak{S}_{E}\) 为支集为 E 的一个轮换， \(\tau\) 为一个对换。证明 \(\mathfrak{S}_{E}\) 由 \(\zeta\) 和 \(\tau\) 生成。
\end{enumerate}

¶ 12. (a) 证明每个置换 \(\sigma \in \mathfrak{A}_n\) 都是 3 阶轮换的乘积（这些轮换一般并不是 \(\sigma\) 的分支轮换）。（先对两个对换的乘积证明之，并利用 § 5 第 7 号命题 9。）

\begin{enumerate}
\def\labelenumi{(\alph{enumi})}
\setcounter{enumi}{1}
\item
  若 \(a_1, \ldots, a_p\) 是 \(p\) 个属于 \([1, n]\) 的互异元素，设 \((a_1 a_2 \ldots a_p)\) 表示 \(p\) 阶轮换，其支集为 \(\{a_1, a_2, \ldots, a_p\}\)，且它把 \(a_i\) 映为 \(a_{i+1}\) （对于 \(1 \leqslant i \leqslant p-1\)），把 \(a_p\) 映为 \(a_1\)。证明 \(\mathfrak{A}_n\) 由 \(n-2\) 个轮换 (1 2 3), (1 2 4), \ldots, (1 2 n) 生成。（利用 (a)。）
\item
  由此推出，如果 \(n\) 是奇数， \(\mathfrak{A}_n\) 由 (1 2 3) 和 (1 2 \ldots{} \(n\) ) 生成，并且，如果 \(n\) 是偶数，则由 (1 2 3) 和 (2 3 \ldots{} \(n\) ) 生成。
\item
  证明：若 \(\mathfrak{A}_n\) 的一个正规子群包含一个 3 阶轮换，则它等于 \(\mathfrak{A}_n\) （证明这样的子群包含所有形如 (12 \(k\) ) 的轮换， \(3 \leqslant k \leqslant n\)）。
\end{enumerate}

¶ 13. 设 G 是集合 X 的一个传递置换群，H 是某元素 \(x \in X\) 的稳定化子。证明以下性质的等价性：

\begin{enumerate}
\def\labelenumi{(\alph{enumi})}
\item
  每个子群 \(\mathbf{H}'\) （\(\mathbf{G}\) 中包含 \(\mathbf{H}\) 者）都等于 \(\mathbf{H}\) 或 \(\mathbf{G}\)。
\item
  X 的每个子集 Y，若对所有 \(g \in G\)，gY 要么包含于 Y，要么与 Y 不相交，则 Y 要么等于 X，要么由单个元素组成。
\end{enumerate}

（如果 \(\mathbf{H}'\) 满足 (a) 且 \(\mathbf{Y} = \mathbf{H}'x\)，则 \(g\mathbf{Y} = \mathbf{Y}\) 对所有 \(g \in \mathbf{H}'\) 成立，且 \(g\mathbf{Y} \cap \mathbf{Y} = \varnothing\) 对所有 \(g \notin \mathbf{H}'\) 成立。反之，若 \(\mathbf{Y}\) 具有性质 (b)，则集合 \(\mathbf{H}'\)（由 \(g \in \mathbf{G}\) 中满足 \(g\mathbf{Y} = \mathbf{Y}\) 的那些 g 组成）是 \(\mathbf{G}\) 的包含 \(\mathbf{H}\) 的一个子群。）

满足条件(a)和(b)且并非仅由恒等元组成的传递置换群 G 称为本原群。

\begin{enumerate}
\def\labelenumi{\arabic{enumi}。}
\setcounter{enumi}{13}
\tightlist
\item
  一个置换群 \(\Gamma\)，作用于集合 \(E\) 上，称为 \(r\) -重传递的，如果对任意两个序列 \((a_1, a_2, \ldots, a_r)\)，\((b_1, b_2, \ldots, b_r)\) （由 \(r\) 个属于 \(E\) 的互异元素组成），存在一个置换 \(\sigma \in \Gamma\) 使得 \(\sigma(a_i) = b_i\) 对 \(1 \leqslant i \leqslant r\) 成立，且该性质对至少一对由 \(r + 1\) 个属于 \(E\) 的互异元素组成的序列不成立。
\end{enumerate}

\begin{enumerate}
\def\labelenumi{(\alph{enumi})}
\item
  证明 \(r\) -重传递群当 \(r > 1\) 时是本原的。
\item
  置换群 \(\Gamma\) 作用于 \(n\) 个对象且为 \(r\) -重传递，其阶具有形式 \(n(n - 1)\ldots (n - r + 1)d\)，其中 \(d\) 是 \((n - r)!\) 的一个因子（考虑 \(\Gamma\) 中保持 \(r\) 个元素不变的置换所成的子群，并计算其指数）。
\end{enumerate}

¶ 15. 设 \(\Gamma\) 是 \(r\) -重传递的置换群，作用于集合 \(\mathbf{E}\)（具有 \(n\) 个元素）上；对于不同于恒等置换的置换 \(\sigma \in \Gamma\)，设 \(n - s\) 为 \(\mathbf{E}\) 中在 \(\sigma\) 下不变的元素的个数。若 \(s > r\)，证明存在一个置换 \(\tau \in \Gamma\) 使得 \(\sigma^{-1}\tau \sigma \tau^{-1}\) 不同于恒等置换，且它在 \(\mathbf{E}\) 中保持不变的元素的个数 \(\geqslant n - 2(s - r + 1)\)（利用 \(\sigma\) 分解为其分支轮换）。如果 \(s = r\)，类似地证明存在 \(\tau \in \Gamma\) 使得 \(\sigma^{-1}\tau \sigma \tau^{-1}\) 是一个 3 阶轮换。由此推出，如果 \(r \geqslant 3\) 且 \(\Gamma\) 不包含交错群 \(\mathfrak{A}_n\)，则 \(s \geqslant 2r - 2\) 对 \(\Gamma\) 中每一个置换成立（利用习题 10）。最后得出结论：如果 \(\Gamma\) 不等于 \(\mathfrak{A}_n\) 或 \(\mathfrak{S}_n\)，则 \(r \leqslant n/3 + 1\)。

¶ 16. (a) 证明交错群 \(\mathfrak{A}_n\) 是 \((n - 2)\) -重传递的。

\begin{enumerate}
\def\labelenumi{(\alph{enumi})}
\setcounter{enumi}{1}
\item
  证明 \(\mathfrak{A}_n\) 对 \(n \geqslant 5\) 是一个非交换单群。（利用 (a)、习题 15 的方法以及习题 12(d)证明 \(\mathfrak{A}_n\) 当 \(n > 6\) 时是单群；类似地考察 \(n = 5\) 和 \(n = 6\) 的情形。）
\item
  证明，对 \(n \geqslant 5\)，\(\mathfrak{S}_n\) 的正规子群只有 \(\{e\}\)，\(\mathfrak{A}_n\) 和 \(\mathfrak{S}_n\)。
\end{enumerate}

\begin{enumerate}
\def\labelenumi{\arabic{enumi}。}
\setcounter{enumi}{16}
\item
  设 G 是集合 X 的一个传递置换群；假设 G 是本原的（习题 13）。证明 G 的每个不同于 \(\{e\}\) 的正规子群都是传递的。
\item
  设 \(\Gamma\) 是集合 E 的一个置换群。设 A 和 B 是 E 的两个子集，互为补集且在 \(\Gamma\) 下稳定。记 \(\Gamma_{A}\) 和 \(\Gamma_{B}\) 为由 \(\Gamma\) 中的置换分别限制到 A 和 B 上所得的群，并记 \(\Delta_{A}\) 和 \(\Delta_{B}\) 为 \(\Gamma\) 中分别保持 A 的每个元素和 B 的每个元素不变的子群。证明 \(\Delta_{A}\) 和 \(\Delta_{B}\) 是 \(\Gamma\) 的正规子群，且 \(\Gamma_{A}\) 同构于 \(\Gamma/\Delta_{A}\)，\(\Gamma_{B}\) 同构于 \(\Gamma/\Delta_{B}\)；若 \(\Delta_{AB}\) （相应地 \(\Delta_{BA}\)）表示 \(\Delta_{A}\) （相应地 \(\Delta_{B}\)）中置换限制到 B（相应地 A）上所得的群，证明商群 \(\Gamma_{A}/\Delta_{BA}\)，\(\Gamma_{B}/\Delta_{AB}\) 和 \(\Gamma/(\Delta_{A}\Delta_{B})\) 是同构的（使用习题 8）。
\item
  设 G 是一个群，H 是 G 的子群，且 G/H 是 G 中（模 H 的）左陪集的集合。设 \(r: \mathrm{G} / \mathrm{H} \to \mathrm{G}\) 是与典范投影 \(\mathrm{G} \to \mathrm{G} / \mathrm{H}\) 相关联的一个截面。则对所有 \(\mathbf{X} \in \mathrm{G} / \mathrm{H}\)，有 \(\mathbf{X} = r(\mathbf{X})\)（模 H）。在 G/H 上通过令 \(\mathbf{X} \top \mathbf{Y} = r(\mathbf{X})r(\mathbf{Y})\)（模 H）定义一个内部合成律。
\end{enumerate}

\begin{enumerate}
\def\labelenumi{(\alph{enumi})}
\item
  证明对所有 X 有 \(X \top H = X\)，且在律 \(\top\) 下每个左平移都是双射。若 \(G'\) 是 G 的由诸元素 \(r(X)\) 组成的集合所生成的子群，且 \(H' = H \cap G'\)，则由映射 r 在 \(G'/H'\) 上类似地定义的内部合成律在该集合上确定了一个结构，它与 G/H 上由律 \(\top\) 所确定的结构同构。
\item
  为使律 \(\top\) 是结合的，必要且充分的是 \(\mathbf{H}'\) 为 \(\mathbf{G}'\) 的正规子群；这时由 \(\top\) 确定的结构同构于商群 \(\mathbf{G}' / \mathbf{H}'\) 的结构（为看出该条件是充分的，首先借助 §4 习题 2(a) 证明：若 \(\top\) 是结合的，则它在 \(\mathbf{G} / \mathbf{H}\) 上确定一个群结构；记 \(\mathbf{K}\) 为 \(\mathbf{G}'\) 的包含于 \(\mathbf{H}'\) 中的最大正规子群，再利用 \(\top\) 的结合性条件证明 \((r(\mathbf{X} \top \mathbf{Y}))^{-1} r(\mathbf{X}) r(\mathbf{Y}) \in \mathbf{K}\) 对所有 \(\mathbf{X}, \mathbf{Y}\) 成立；由此推出映射 \(\mathbf{X} \mapsto r(\mathbf{X})\)（模 \(\mathbf{K}\)）是群 \(\mathbf{G} / \mathbf{H}\)（带律 \(\top\)）到商群 \(\mathbf{G}' / \mathbf{K}\) 上的一个同构；由此得出 \(\mathbf{H}' = \mathbf{K}\)，为此注意 \(\mathbf{H}'\) 是模 \(\mathbf{K}\) 的若干陪集之并。）
\item
  反之，假设在集合 E 上给定一个内部合成律 \(\top\)，使得每个左平移都是双射，且存在 \(e \in \mathbf{E}\) 使 \(x \top e = x\) 对所有 \(x \in \mathbf{E}\) 成立。设 \(\Gamma\) 为 E 的由诸左平移 \(\gamma_x\) 生成的置换群，\(\Delta\) 为 \(\Gamma\) 中使 \(e\) 不变的置换所成的子群；证明对模 \(\Delta\) 的每一个左陪集 X，恰好有一个元素 \(x \in \mathbf{E}\) 使得 \(\gamma_x \in \mathbf{X}\)；若记 \(r(\mathbf{X}) = \gamma_x\)，则映射 \(x \mapsto \gamma_x\) 是带律 \(\top\) 的集合 E 到集合 \(\Gamma / \Delta\) （带律 \((\mathbf{X}, \mathbf{Y}) \mapsto r(\mathbf{X})r(\mathbf{Y})\)（模 \(\Delta\)））上的一个同构。
\end{enumerate}

\begin{enumerate}
\def\labelenumi{\arabic{enumi}。}
\setcounter{enumi}{19}
\tightlist
\item
  设 \(G\) 是一个单群， \(H\) 是 \(G\) 的一个指数为有限数 \(n > 1\) 的子群。
\end{enumerate}

\begin{enumerate}
\def\labelenumi{(\alph{enumi})}
\item
  证明 \(\mathbf{G}\) 是有限的，且其阶整除 \(n!\) （使用习题 5）。
\item
  证明 \(\mathbf{G}\) 当 \(n \leqslant 4\) 时是交换的（使用习题 10）。
\end{enumerate}

\begin{enumerate}
\def\labelenumi{\arabic{enumi}。}
\setcounter{enumi}{20}
\tightlist
\item
  设 \(p(n)\) 为对称群 \(\mathfrak{S}_n\) 的共轭类个数。证明 \(p(n)\) 等于族 \((x_1, \ldots, x_n)\)（由整数 \(\geqslant 0\) 组成）的个数，这些族满足 \(\sum_{i=1}^{n} i.x_i = n\)。
\end{enumerate}

* 推导恒等式 \(\sum_{n=0}^{\infty} p(n) T^n = \prod_{m=1}^{\infty} \frac{1}{1 - T^m}\) .*

\begin{enumerate}
\def\labelenumi{\arabic{enumi}。}
\setcounter{enumi}{21}
\tightlist
\item
  设 \(A = Z/2Z\)，\(B = S_{3}\) 和 \(G = A \times B\)。若 s 是 B 中阶为 2 的元素，令 \(\phi\) 为 G 的自同态，其核等于 B，其像为 \(\{1, s\}\)。证明 G 的中心在 \(\phi\) 下不是稳定的。
\end{enumerate}

¶ 23. (a) 设 s 是群 \(\mathfrak{S}_n\) 中阶为 2 的一个元素，它是 \(k\) 个支集两两不相交的对换 \(t_1, \ldots, t_k\) 的乘积。证明中心化子 \(\mathrm{C}_s\) （\(s\) 在 \(\mathfrak{S}_n\) 中的）同构于 \(\mathfrak{S}_{n-2k} \times \mathrm{A}\)，其中 A 有一个 \(2^k\) 阶的正规子群 B，使得 A/B 同构于 \(\mathfrak{S}_k\) （注意， \(\mathrm{C}_s\) 的每个元素都置换 \(s\) 的不动点以及 \(t_1, \ldots, t_k\) 的支集）。

\begin{enumerate}
\def\labelenumi{(\alph{enumi})}
\setcounter{enumi}{1}
\item
  证明，如果 \(n = 5\) 或 \(n \geqslant 7\)，\(\mathfrak{S}_n\) 中每个中心化子与一个对换的中心化子同构的 2 阶元素本身是一个对换（利用习题 16 比较这些中心化子的若尔当–赫尔德商）。由此推出， \(\mathfrak{S}_n\) 的每个自同构都置换诸对换。
\item
  若 \(a, b\) 是 \([1, n]\) 的两个互异元素，设 \(F_a\) （相应地 \(F_{ab}\) ) 表示 \(a\) 的稳定化子（相应地， \(\{a, b\}\) 的固定子）在 \(\mathfrak{S}_n\) 中，并设 \((ab)\) 为对换 \(\tau_{a,b}\)。设 \(a, b, c, d\) 是 \([1, n]\) 的互异元素。证明，如果 \(n \geqslant 5\)，由 \(F_{ab}\) 和 \(F_{cd}\) 生成的群是 \(\mathfrak{S}_n\) （证明这个群包含所有对换；为此使用如下事实：如果 \(x \in [1, n] = \{a, b, c, d\}\)，则 \((ac) = (cx)(ax)(cx)\)）。证明由 \(F_{ab}\) 和 \(F_{ac}\) 生成的群是 \(F_a\) （使用与上述情形相同的等式）。
\item
  证明一个自同构 \(\alpha\) （\(\mathfrak{S}_n\) 上的）若满足 \(\alpha(F_a) = F_a\) 对所有 \(a \in [1, n]\)，便是恒等映射（证明 \(\alpha\) 把每个对换映射到自身）。证明，如果存在 \(a \in [1, n]\) 使得 \(\alpha(F_a) = F_a\)，那么 \(\alpha\) 是一个内自同构（证明存在一个内自同构 \(\beta\) 使得 \(\beta \circ \alpha\) 把每个群 \(F_b\)，\(b \in [1, n]\) 映射到自身）。
\item
  设 \(a, b\) 是 \([1, n]\) 的两个互异元素，并设 \(\mathbf{C}_{ab}\) 是 \((a, b)\) 在 \(\mathfrak{S}_n\) 中的中心化子。证明，如果 \(n \geqslant 5\)，\(\mathbf{C}_{ab}\) 的导群在 \(\mathbf{C}_{ab}\) 中的指数为 4（利用如下事实： \(\mathfrak{S}_{n-2}\) 的导群是 \(\mathfrak{A}_{n-2}\)）。证明 \(\mathbf{F}_{ab}\) 是 \(\mathbf{C}_{ab}\) 的唯一一个指数为 2、不包含 \((a, b)\) 且不包含于 \(\mathfrak{A}_n\) 的子群。
\item
  证明，如果 \(n \geqslant 5\)，自同构 \(\alpha\) （\(\mathfrak{S}_n\) 的）若保持一个对换 \((ab)\) 不变，便是内自同构。（借助 (e) 证明 \(\alpha(\mathrm{F}_{ab}) = \mathrm{F}_{ab}\)；同样证明，如果 \(c \notin \{a, b\}\)，\(\alpha(\mathrm{F}_{ac})\) 或者形如 \(F_{ax}, x \notin \{a, b\}\)，或者形如 \(\mathrm{F}_{bx}, x \notin \{a, b\}\)；在第一种情形，得出 \(\alpha(\mathrm{F}_a) = \mathrm{F}_a\) 并应用 (d)；在第二种情形，将 \(\alpha\) 乘以 \(\mathfrak{S}_n\) 的由 \((ab)\) 定义的内自同构，把它归结为第一种情形。）
\item
  从 (b) 和 (f) 推出， \(\mathfrak{S}_n\) 的每个自同构当 \(n = 5\) 或 \(n \geqslant 7\) 时是内自同构。
\end{enumerate}

\begin{enumerate}
\def\labelenumi{\arabic{enumi}。}
\setcounter{enumi}{23}
\tightlist
\item
  \begin{enumerate}
  \def\labelenumii{(\alph{enumii})}
  \tightlist
  \item
    证明存在 \(\mathfrak{S}_6\) 的一个子群 H，它同构于 \(\mathfrak{S}_5\) 且不固定 \([1,6]\) 的任何元素（令 \(\mathfrak{S}_5\) 通过内自同构作用于其 5 阶子群的集合，该集合有 6 个元素）。
  \end{enumerate}
\end{enumerate}

\begin{enumerate}
\def\labelenumi{(\alph{enumi})}
\setcounter{enumi}{1}
\item
  证明存在一个自同构 \(\sigma\) （\(\mathfrak{S}_{6}\) 的），它把 (12) 映为 (12)(34)(56)。（令 \(\mathfrak{S}_{6}\) 作用于 \(\mathfrak{S}_{6} / \mathrm{H}\)，其中 \(\mathrm{H}\) 取为如上所选取者。）
\item
  证明 \(\mathfrak{S}_{6}\) 的内自同构群在 \(\mathfrak{S}_{6}\) 的全自同构群中的指数为 2。（设 \(\alpha\) 是 \(\mathfrak{S}_{6}\) 的一个自同构；证明 \(\alpha\) 把一个对换要么映为一个对换，要么映为 (1 2)(3 4)(5 6) 的一个共轭。利用习题 23，证明在第一种（相应地，第二种）情形下， \(\alpha\) （相应地 \(\alpha \circ \sigma\) )是内自同构。）
\end{enumerate}

\begin{enumerate}
\def\labelenumi{\arabic{enumi}。}
\setcounter{enumi}{24}
\tightlist
\item
  设 \(n\) 为整数 \(\geqslant 5\)。证明 \(\mathfrak{S}_n\) 的阶为 \((n - 1)!\) 的子群都同构于 \(\mathfrak{S}_{n-1}\)，并且当 \(n \neq 6\) （相应地当 \(n = 6\)）时构成单一的共轭类（相应地两个共轭类）。（设 \(H\) 为这样的一个子群， \(x_1, \ldots, x_n\) 为 \(\mathfrak{S}_n / \mathrm{H}\) 的元素。证明 \(\mathfrak{S}_n\) 在诸 \(x_i\) 上的作用定义了 \(\mathfrak{S}_n\) 的一个自同构，并应用习题 23 和 24。）
\end{enumerate}

¶ 26. (a) 设 G 是作用于两个有限集合 \(E_{1}\) 和 \(E_{2}\) 上的一个群。若 \(s \in G\)，设 \(s_{1}\) （相应地 \(s_{2}\) ) 表示由 s 在 \(E_{1}\) （相应地 \(E_{2}\) ) 上定义的置换，并设 \(E_{1}^{s}\) （相应地 \(E_{2}^{s}\) ) 为在 \(s_{1}\) （相应地 \(s_{2}\) ) 下不变的元素组成的集合。

证明以下性质的等价性：

\begin{enumerate}
\def\labelenumi{(\roman{enumi})}
\item
  对所有 \(s \in G\)，有 \(\operatorname{Card}(\mathbf{E}_1^s) = \operatorname{Card}(\mathbf{E}_2^s)\)。
\item
  对所有 \(s \in \mathbf{G}\)， \(s_1\) 的分支轮换的阶（参见第 7 号命题 7）与 \(s_2\) 的分支轮换的阶（至多差一个置换）相同。
\item
  对所有 \(s \in \mathbf{G}\)，存在一个双射 \(f_s: \mathbf{E}_1 \to \mathbf{E}_2\) 使得 \(s_2 \circ f_s = f_s \circ s_1\)。
\end{enumerate}

如果这些性质成立，则称 G-集合 \(E_{1}\) 和 \(E_{2}\) 是弱等价的。

\begin{enumerate}
\def\labelenumi{(\alph{enumi})}
\setcounter{enumi}{1}
\item
  给出两个弱等价但不同构的 G-集合的例子（取 G 为一个四阶非循环群）。
\item
  设 \(H_{1}\) 和 \(H_{2}\) 是有限群 G 的两个子群。证明以下性质的等价性：
\end{enumerate}

\begin{enumerate}
\def\labelenumi{(\arabic{enumi})}
\tightlist
\item
  对于 G 的每个共轭类 C，
\end{enumerate}

\[
\operatorname{Card} (\mathrm{C} \cap \mathrm{H} _ {1}) = \operatorname{Card} (\mathrm{C} \cap \mathrm{H} _ {2}).
\]

\begin{enumerate}
\def\labelenumi{(\arabic{enumi})}
\setcounter{enumi}{1}
\tightlist
\item
  G-集 \(\mathbf{G} / \mathbf{H}_1\) 和 \(\mathbf{G} / \mathbf{H}_2\) 是弱等价的。
\end{enumerate}

\begin{enumerate}
\def\labelenumi{(\alph{enumi})}
\setcounter{enumi}{3}
\item
  证明，若 \(H_{1}\) 和 \(H_{2}\) 满足上述性质 (1) 和 (2)，且 \(H_{1}\) 在 G 中正规，则 \(H_{2} = H_{1}\)。
\item
  设 \(\mathbf{G} = \mathfrak{A}_6\)，并设
\end{enumerate}

\[
\begin{array}{l} \mathbf {H} _ {1} = \{e, (1 2) (3 4), (1 3) (2 4), (1 4) (2 3) \} \\ \mathbf {H} _ {2} = \{e, (1 2) (3 4), (1 2) (5 6), (3 4) (5 6) \}. \end{array}
\]

证明 \(H_{1}\) 和 \(H_{2}\) 满足性质 (1) 和 (2)，并且 \(H_{1}\) 和 \(H_{2}\) 在 \(A_{6}\) 中不共轭（在 \(S_{6}\) 中也不共轭）。

\begin{enumerate}
\def\labelenumi{\arabic{enumi}。}
\setcounter{enumi}{26}
\item
  设 Y 是集合 X 的一个子集，并设 A 是 X 的置换群 \(S_{X}\) 中 Y 的固定子。设 M 为 \(S_{X}\) 的由满足 \(sAs^{-1} \subset A\) 的元素 s 组成的子幺半群。证明 M 是 \(S_{X}\) 的子群当且仅当集合 Y 与 X = Y 中有一个是有限的。
\item
  设 G 是一个群，A 和 B 是 G 的两个子群，且 \(\phi\) 是 A 到 B 上的一个同构。证明存在一个包含 G 的群 \(G_{1}\)，使得 \(\phi\) 是 \(G_{1}\) 的一个内自同构的限制（使用集合 G 的置换群）。证明：若 G 是有限的，则 \(G_{1}\) 可选取为有限的。
\item
  设 X 为无限集。设 G 为 \(S_{X}\) 的由具有以下性质的置换 \(\sigma\) 组成的子群：存在 X 的一个有限子集 \(Y_{\sigma}\)，使得 \(\sigma x = x\) 对 \(x \in X - Y_{\sigma}\) 成立，且 \(\sigma\) 在 \(Y_{\sigma}\) 上的限制是偶的。
\end{enumerate}

\begin{enumerate}
\def\labelenumi{(\alph{enumi})}
\item
  证明 G 是一个非交换单群。证明 G 的每个生成 G 的子集都与 X 等势。
\item
  证明每个有限群都同构于 G 的一个子群（利用以下事实：\(S_{n}\) 同构于 \(A_{n+2}\) 的一个子群）。
\end{enumerate}

